\subsection{Efficient focusing with curved crystals: Johann geometry and Rowland circle}

Consider a crystal bent to radius $R$, with diffraction planes \textbf{parallel to the
crystal surface} (symmetric cut, asymmetry angle $\chi = 0$).  A point source $S$ at
distance $p$ from the crystal centre emits a divergent beam that impinges on the concave
surface.  The condition on $p$ and the image distance $q$ such that the reflected
rays converge to a focus is given by the Coddington equation (for the meridional or diffraction plane)


\begin{equation}
  \frac{1}{p} + \frac{1}{q} = \frac{2}{R\sin\theta_B}
  \label{eq:meridional_sym}
\end{equation}
where
$\\theta_B$ is the Bragg angle at the crystal centre (it is focused independently in the sagittal plane at a different distance, but this is not discussed here).

The \textbf{Johann geometry} bends a flat crystal to radius $R$ and places source $S$ and
detector $D$ symmetrically at the \emph{equal-conjugate} distance (i.e. 1:1 magnification) satisfying Eq.~\eqref{eq:meridional_sym}:
\begin{equation}
  p = q = R\sin\theta_B.
  \label{eq:equal_conj}
\end{equation}

At the crystal \emph{centre}, this placement is exact: the incidence angle equals $\\theta_B$
and the Bragg condition $\lambda = 2d\sin\\theta_B$ is perfectly satisfied.

Origin of the Johann Error: 
Now consider a ray hitting the crystal at a point $P$ displaced by a distance $\ell$ along
the crystal surface from the centre.  Because the crystal is \emph{cylindrically} bent, its
surface follows a circle of radius $R$.  The local surface normal at $P$ is rotated by an
angle
\begin{equation}
  \delta\phi = \frac{\ell}{R}
  \label{eq:normal_rotation}
\end{equation}
relative to the central normal.  The source $S$, however, is still at the fixed position
defined by Eq.~\eqref{eq:equal_conj}.  The actual incidence angle at $P$ therefore deviates
from $\\theta_B$, and the Bragg condition is no longer perfectly satisfied.

Working out the geometry to leading order in $\ell/R$, the \textbf{angular error} at $P$ is:
\begin{equation}
  \delta\theta \approx \frac{\ell^2 \cos^2\theta_B}{2R^2\sin\theta_B}.
  \label{eq:angular_error}
\end{equation}

This error is \emph{second order} in $\ell$: it vanishes at the crystal centre and grows
towards the edges.  Via the differential of Bragg's law
$\delta E / E = -\cot\\theta_B\,\delta\theta$, the corresponding relative energy broadening
is:

\begin{keyresult}
\begin{equation}
  \delta E \;\approx\; \frac{\ell^2\cos^2\theta_B}{2R^2}
  \label{eq:johann_error}
\end{equation}
\end{keyresult}

This is the \textbf{Johann error}.  Its salient properties are:
\begin{itemize}
  \item It scales as $\ell^2$, so using a smaller crystal reduces the error quadratically,
        at the cost of collected solid angle.
  \item It vanishes at $\\theta_B = 90°$ (backscattering), since $\cos 90° = 0$.
  \item For a given crystal size $\ell$, larger bending radius $R$ reduces the error ---
        but also reduces the collected solid angle.
  \item It is an \emph{intrinsic geometrical aberration} of cylindrical bending: the
        crystal surface cannot conform simultaneously to the Bragg condition at every point.
\end{itemize}

The Rowland Circle Construction.
The three defining elements are:
\begin{itemize}
  \item \textbf{Source $S$}, \textbf{crystal centre $C$}, and \textbf{detector $D$} all
        lie on the Rowland circle of radius $R_W = R/2$.
  \item The crystal is tangent to the Rowland circle at $C$, and bent to radius
        $R = 2R_W$.
  \item The chord $SD$ subtends an angle $2\\theta_B$ at the centre of the Rowland circle,
        consistent with Bragg's law.
\end{itemize}

As energy (and hence $\\theta_B$) changes, the focus point $D$ moves continuously along the
Rowland circle.  This is the basis for \textbf{energy scanning in Johann spectrometers}
without translating the source.


\paragraph{Extension to the Asymmetric Case}
In the \textbf{asymmetric Bragg geometry}, the diffraction planes make an angle $\chi$
(the \emph{asymmetry angle}) with the crystal surface.  The incidence and exit angles with
respect to the \emph{surface} are then:
\begin{equation}
  \alpha = \theta_B - \chi, \qquad \beta = \theta_B + \chi.
  \label{eq:asym_angles}
\end{equation}
(Signs are for the beam-compressing case; they reverse for beam expansion.)  Both $\alpha$
and $\beta$ satisfy the Bragg condition with respect to the diffracting planes, since the
\emph{planes} are still at angle $\\theta_B$ to the incident beam.

The asymmetry introduces a \textbf{beam expansion factor}:
\begin{equation}
  b = -\frac{\sin\alpha}{\sin\beta} = -\frac{\sin(\theta_B-\chi)}{\sin(\theta_B+\chi)}.
  \label{eq:asym_ratio}
\end{equation}
Here $|b|<1$ corresponds to angular expansion (spatial compression) and $|b|>1$ to angular
compression (spatial expansion).


Applying Fermat's principle to the asymmetric geometry --- where the incidence and exit
angles on the crystal surface are $\alpha$ and $\beta$ respectively --- modifies the
Coddington equation:

\begin{keyresult}
\begin{equation}
  \frac{\sin^2\alpha}{p} + \frac{\sin^2\beta}{q}
  = \frac{\sin\alpha + \sin\beta}{R}
  \label{eq:meridional_asym}
\end{equation}
\end{keyresult}

For perfect meridional focusing from a point source, the equal-conjugate condition
generalises to:
\begin{equation}
  p = \frac{2R\sin^2\alpha}{\sin\alpha+\sin\beta},
  \qquad
  q = \frac{2R\sin^2\beta}{\sin\alpha+\sin\beta}.
  \label{eq:asym_conjugates}
\end{equation}

The source and detector are now placed at \emph{different} distances from the crystal, but
both still lie on the Rowland circle.  To see this, note that the Rowland circle radius is
unchanged:

\begin{keyresult}
\begin{equation}
  R_W = \frac{R}{2}
  \label{eq:rowland_asym}
\end{equation}
\end{keyresult}

The Rowland circle radius is \textbf{independent of asymmetry}.  What changes is the
angular position of $S$ and $D$ on the circle: $S$ sits at arc angle $\alpha$ and $D$ at
arc angle $\beta$ from the crystal normal, so the crystal is no longer at the midpoint of
the $SD$ arc.  The inscribed angle theorem still guarantees that the local Bragg condition
is satisfied at every surface point that lies on the Rowland circle (Johansson case), or
approximately so for the Johann approximation.
